\chapter{C14.3b: a complete source subchart from a damped two-circuit completion}
\label{component:p024c821bfb59}
\noindent\textbf{Component }\path{p024c821bfb59}\par
{\small\noindent Source: \path{R17/proofs/c1_small_eta/damped_gamma_source_subchart.tex}\par}
{\small\noindent SHA256: \path{024c821bfb59309356b1562279a3ead111f5821540325d92c293cf2376932597}\par}
{\small\noindent Inventory use: DG2--DG9 and all allocations except DG10 | C14\_\allowbreak{}ENTIRE\_\allowbreak{}ORIGINAL\_\allowbreak{}PARENT\_\allowbreak{}GRAPH; precise subordinate uses retained below\par}
\medskip
\section{C14.3b: a complete source subchart from a damped two-circuit completion}

\begin{theorem}
Normalize $d=1$ and retain the unchanged physical C14.3b source, with $x,\eta>0$
infinitesimal, $\zeta$ a positive bounded unit, every original $j$
and $1-r$ opening, $\beta\ge0$ including literal $\beta=0$, and all inherited intersections.
On the additional \emph{pure-source} subchart
\[
                         x/\eta\text{ bounded},              \tag{DG1}
\]
the entire original parent is finite tensor SOS for every independent
projective angle and state test. No relation between a source rate
and a testing rate is assumed. The proof below treats the whole local
corner; the inherited remote theorem gives its complement. The source
subchart $x/\eta$ unbounded is not admitted by this theorem.
\end{theorem}

\paragraph{The sign reduction and the complete frame.}
Use the sign gauges to take $u,v\ge0$. For real states with $WZ\le0$
the independent half-core sign lemma proved in this revision gives a
direct tensor-SOS decomposition of the original signed half core;
all remaining original event packets are positive source multiples
of squares. Thus only $WZ>0$ needs the following construction. No
normalization or inverse of a state coordinate is used here.

Use all homogeneous JF1 definitions, and abbreviate $Q=Q_t$,
$T=T_D$, $B=B_x$, $a=1/\zeta$, $x=k^2$, $a=\eta c^2$. In particular
\[
 E=\eta uW-vZ=\eta(P+Q),\quad H=avW-uZ,
 \quad \tau=\eta cW-Z,\quad D_0=u-cv,
\]
\[
 J=(uv-x/c)Z,\quad
 R=-\sqrt\lambda\,\eta kD_0W,\quad
 f_t=F_t/\eta,\quad \mathcal Q=\sqrt q(H-a\xi),
 \quad N_x=\sqrt B\xi.
\]
The full radial packet $\mathscr R$ is JF2, with the strict two-root
completion JF3. Its roots will be denoted $R_b,f_t$; both normalized
pivots have positive unit residue. The angular packet is exactly
\[
 \mathscr A=\mathcal Q^2
       -\frac{2\mu}{a\sqrt{qB}}\mathcal QN_x
       +(1-\mu/B)N_x^2.                              \tag{DG2}
\]
Its normalized Schur pivot is strict by the unchanged R15/R16 source
margin, including $\beta=0$ and unbounded $B$. Those previously
proved angular and radial blocks are used once, as their complete
packets.

\paragraph{Damp the remaining normal and complete both circuit companions.}
Set the positive source unit
\[
 G_0=2\eta\lambda-2\eta\sim1,\qquad
 \epsilon=\gamma-G_0/h,
\]
\[
 \widehat R=R+\frac{G_0k}{2\sqrt\lambda}Q,
 \qquad \widehat J=\sqrt\lambda J+
                         \frac{G_0u}{2\sqrt\lambda}Q,
 \qquad c_U=A-\frac{G_0^2}{4\lambda}.
\]
Use the exact full-$D$ conductor and polarization GS1 to conduct only
$(G_0/h)uvQ\xi$, leaving $\epsilon uvQ\xi$ unaltered. Complete
\emph{both} $R$ and $J$, incorporating the entire native $AUQ^2$
row into the latter completion. Since
$G_0/(2\lambda)=\eta-\eta/\lambda$, the complete original parent
has the following exact identity:
\[
\begin{split}
 F/4={}&\mathscr A+\mathscr R+S_0
    -\lambda x(\eta P+\eta Q/\lambda)^2
    +\widehat R^2+\widehat J^2\\
 &+c_UUQ^2+e_4UF_t^2+BV\xi^2
    +\frac{2x}{c}uv\tau^2
    +\frac{G_0x}{c}D_0^2WZ\\
 &-e_4yuHF_t+\kappa uF_t\xi
                 +\frac{G_0x}{c}PH+\epsilon uvQ\xi,       
\end{split}\tag{DG3}
\]
where
\[
 S_0=CP^2+AQ^2+D_p(P-hT)^2+2b_0QT,
 \qquad b_0=b'-G_0/2.                                 \tag{DG4}
\]
This is a simultaneous completion of the entire remaining normal and
the coupled circuit. It retains every full-$D$, full-$Y$ and higher
normal of JF1. There is no second bare radial pivot and no extracted
uncentered $m$ event. The coefficient of the extra $uv\tau^2$ row
is exactly $2x/c>0$, because $2\lambda-G_0/\eta=2$.

\paragraph{Two strictly positive source deficits replace the failed soft pivot.}
Put $K=CD_p/(C+D_p)$, $s=x+\eta$, and
\[
                         \Delta_0=h^2AK-b_0^2.
\]
The original source formulas imply exact finite divisions
\[
 \Delta_0=x d_x+\eta d_\eta,\qquad
 \frac{c_U}{\eta}=x e_x+\eta e_\eta,                \tag{DG5}
\]
with bounded source coefficients having the following residues:
\[
 \begin{array}{c|c}
 d_x&\pi/2+1/4+1/(4a)\\
 d_\eta&\sigma/2+3/4\\
 e_x&\pi+1+1/(2a)\\
 e_\eta&\sigma+3/2
 \end{array}
 \qquad h=1-x\pi,\quad r=1-\eta\sigma.             \tag{DG6}
\]
All four residues are strictly positive. Here $\pi,\sigma\ge0$
are bounded by the unchanged physical source estimates. Thus
\[
 \Delta_0/s\text{ is a positive bounded unit},\qquad
 c_U/(\eta s)\text{ is a positive bounded unit}.     \tag{DG7}
\]
These assertions compare source quantities only and do not use DG1.

For clarity, the divisions in DG5 are opening-safe. Write
$C_n=(1-r)C$, $D_j=jD_p$ and $\Theta_B=B(1-h)$, all given by their
unchanged unit-denominator source formulas. The identity
$\eta\Theta_B-D_j=-C_n$ gives
$b'-\gamma h/2=-hC_n/r$. The gap-cleared harmonic coefficient is
\[
 \eta K=\frac{C_nD_j}
 {C_n(1+xh-\sigma)+D_j\sigma}.
\]
Its denominator is a positive unit, even at either opening. The
quantities $\eta\lambda$, $b'$, $A/\eta$ and $\gamma h$ also have
unit-denominator expressions after their displayed monomial factors
are canceled. Both functions on the left of DG5 vanish at
$x=\eta=0$. Define their $\eta$ coefficients by their exact $x=0$
values divided by $\eta$, and then divide the remainder by $x$.
This gives DG5 with only source-unit denominators; differentiation at
the origin gives DG6. No assumption on $x/\eta$ enters these exact
finite divisions.

\paragraph{The source roots control the complete soft direction.}
There is the exact pure-source completion
\[
 S_0=N_s^2+N_T^2+N_Q^2,
\]
\[
 N_s=\sqrt{C+D_p}\left(P-\frac{hD_p}{C+D_p}T\right),
 \quad N_T=h\sqrt K\left(T+\frac{b_0}{h^2K}Q\right),
 \quad N_Q=\sqrt{\frac{\Delta_0}{h^2K}}Q.            \tag{DG8}
\]
As $\eta K\sim1/2$, DG7 gives
$N_Q/\{\sqrt{\eta s}\,Q\}$ a positive source unit. The inverse
identities show
\[
 T=\frac{N_T}{h\sqrt K}-\frac{b_0}{h^2K}Q,
 \qquad
 P=\frac{N_s}{\sqrt{C+D_p}}
                      +\frac{hD_p}{C+D_p}T.         \tag{DG9}
\]
Consequently $P$ has bounded coefficients on $(N_s,N_T,N_Q)$:
the first two coefficients are $O(\sqrt\eta)$ and the last is
$O(\sqrt{\eta/s})$, which is bounded because $\eta\le s$.
This is the complete $P$ control after the full $D$ shift; it is not
an assumption that the old $E/\sqrt\eta$ pivot remains strict.

The negative square in DG3 is an infinitesimal Gram on these three
roots. Its root is
\[
 \sqrt{\lambda x}(\eta P+\eta Q/\lambda)
   =\eta\sqrt{\lambda x}\,P+\eta\sqrt{x/\lambda}\,Q.
\]
Since $\eta\lambda\sim1/2$, the first coefficient is
$O(\sqrt{\eta x})$ on the bounded allocation of $P$; the second,
after normalizing $Q$ by $N_Q$, is
$O(\eta\sqrt{x/s})$. Both are infinitesimal, without DG1.

\paragraph{Every remaining complete normal is small on this frame.}
The angular-radial term in DG3 is
\[
 -e_4y\eta u\left(\mathcal Q/\sqrt q
                             +aN_x/\sqrt B\right)f_t.
\]
Its coefficients are $O(u)$ and $O(xu)$, using
$e_4\eta^2\sim1/2$ and $q/x^2$ a positive bounded unit.
The term $\kappa uF_t\xi$ has coefficient
$(\eta\kappa)u/\sqrt B$, again infinitesimal.
The new $PH$ term has coefficients
\[
 \frac{G_0x}{c\sqrt q}\quad\hbox{and}\quad
 \frac{G_0ax}{c\sqrt B}
\]
on $P\mathcal Q$ and $PN_x$; these are respectively
$O(\sqrt\eta)$ and $O(x\sqrt\eta)$, so DG9 allocates the whole
term infinitesimally.

Finally $\epsilon=O(s)$ by the same gap-free source expressions.
The complete, still unconducted normal has coefficient
\[
 \frac{\epsilon uv}{\sqrt{\eta s}\sqrt B}
         =O\!\left(\sqrt{s/\eta}\,uv\right)          \tag{DG10}
\]
on $N_QN_x$. DG1 is used \emph{exactly here}: it makes $s/\eta$
bounded, and hence the coefficient infinitesimal for independent
infinitesimal $u,v$. No testing comparison is imposed.

Thus the remainder on
\[
                  (N_s,N_T,N_Q,R_b,f_t,\mathcal Q,N_x)
\]
has a bounded Gram with strictly positive residue. The bounded-Gram
lemma makes it finite tensor SOS. Add the whole positive rows
$\widehat R^2$, $\widehat J^2$, $c_UUQ^2$, $e_4UF_t^2$,
$BV\xi^2$, $(2x/c)uv\tau^2$ and
$(G_0x/c)D_0^2WZ$. The last is a literal tensor square on the current
pure-testing sign chart; all radicals are source-only or testing-only.
This proves the complete local parent, not merely its scalar
restriction.

\paragraph{Equality, endpoints and the global join.}
At original contact $P,Q,T,J,R,\tau,F_t,\xi$ vanish, so all newly
completed rows vanish too. DG3 preserves all contact normals exactly.
The row formulas have no angle ratio or state-coordinate inverse;
$\sqrt{uv}$ is defined as zero on a literal axis. Thus they include
both literal local axes. The inherited gauges cover every angle sign.
For a complex state, split the real-coefficient Hermitian parent into
its real and imaginary evaluations and apply either the $WZ\le0$
core lemma or this $WZ>0$ proof separately to each real state. There
is no assumption that both real and imaginary state components have
a unit $Z$ or $W$. Combine the local proof with the unchanged remote
all-state theorem to obtain the all-tests assertion on DG1.

The focused checker \path{check_damped_gamma_closure.py}
verifies the new two-circuit completion, exact source square roots,
and the four source coefficients in DG6. The companion GS checker
certifies the full affected-block identity. This is a genuine source
subchart admission; it leaves the source chart $x/\eta$ unbounded
open and makes no inference there from DG10.
